\chapter*{Tự đánh giá báo cáo}
\addcontentsline{toc}{chapter}{Tự đánh giá báo cáo}
Nhóm tự đối chiếu báo cáo với yêu cầu nội dung và trình bày của môn học (đánh dấu \checkmark{} khi đạt, ghi chú nếu chưa).

\begin{center}\small\renewcommand{\arraystretch}{1.3}
\begin{tabularx}{\textwidth}{|>{\raggedright\arraybackslash}X|c|>{\raggedright\arraybackslash}p{4cm}|}
\hline
\textbf{Yêu cầu} & \textbf{Đạt} & \textbf{Vị trí / ghi chú}\\ \hline
Phần I: Tổng quan về kho dữ liệu & & Chương 1\\ \hline
Phần II: Tổng quan về BI & & Chương 2\\ \hline
Phần III: Ứng dụng vào bài toán: giới thiệu bài toán (requirement) & & Mục 3.1, 3.2.1--3.2.2\\ \hline
Giới thiệu về ODS & & Mục 3.3.3, 3.4.2\\ \hline
Kiến trúc kho dữ liệu (có logo công nghệ ở mỗi phần) & & Mục 3.3.2\\ \hline
Tiền xử lý dữ liệu (có before--after) & & Mục 3.3.1, 3.4.2--3.4.3\\ \hline
Mô hình dữ liệu OLTP (gạch chân trường dùng sang OLAP, ghi mapping) & & Mục 3.2.4, Phụ lục C\\ \hline
Mô hình dữ liệu OLAP (data model logic, vật lý; fact/dimension khác màu) & & Mục 3.3.4\\ \hline
Các mẫu dashboard (đa dạng biểu đồ, slicer, định dạng có điều kiện) & & Mục 3.4.5\\ \hline
Có mục lục, kết luận, tài liệu tham khảo, có số trang & & \\ \hline
Font Times New Roman 12, căn lề hai bên & & \\ \hline
Không sai chính tả & & \\ \hline
\end{tabularx}
\end{center}
\todo{Nhóm điền cột "Đạt" và ghi chú khi hoàn thiện báo cáo.}
